\let\cleardoublepage\clearpage
\chapter{Горное и нефтегазовое дело}
\ID{МРНТИ 52.01.81, 50.47.31}{}

\begin{header}
\swa{К.~Акишев, А.~Зейнуллин, Е.~Айбульдинов, А.~Подвалов, А.~Нурмаганбетов, М.~Оспанова, М.~Чагай, М.~Айтиева, С.~Хусаинова, Л.~Мылтыкбаева}{СИСТЕМА ОБНАРУЖЕНИЯ И УДАЛЕНИЯ НЕДРОБИМЫХ ВКЛЮЧЕНИЙ НА КОНВЕЙЕРНЫХ ЛИНИЯХ ГОРНО-ОБОГАТИТЕЛЬНЫХ КОМБИНАТОВ}

\tsp{1}К.~Акишев\envelope,
\tsp{2}А.~Подвалов,
\tsp{1}А.~Зейнуллин\envelope,
\tsp{1}Е.~Айбульдинов\envelope,
\tsp{2}А.~Нурмаганбетов,
\tsp{1}М.~Оспанова,
\tsp{2}М.~Чагай,
\tsp{2}М.~Айтиева,
\tsp{3}С.~Хусаинова,
\tsp{1}Л.~Мылтыкбаева
\end{header}

\begin{affil}
\tsp{1}Казахский университет технологии и бизнеса им.К.Кулажанова, Астана, Казахстан,

\tsp{2}ТОО AG Tech, Астана, Казахстан,

\tsp{3}Евразийский национальный университет им. Л.Н. Гумилева, Астана, Казахстан

\corrauthor{Корреспонденты-авторы: akmail04cx@mail.ru, karim\_57@mail.ru}
\end{affil}

Актуальность исследования связана прежде всего с высокой аварийностью и
значительными эксплуатационными потерями конвейерных линий
рудоподготовки горно-обогатительных комбинатов, из-за попадания
недробимых включений в технологическое оборудование. Целью работы
является разработка и внедрение интеллектуальной системы обнаружения и
автоматического удаления недробимых включений для повышения надёжности и
эффективности работы дробильно-транспортного комплекса (ДТК).

В работе использованы методы анализа эксплуатационных данных, алгоритмы
машинного зрения, интеллектуальной обработки сигналов, методы системного
и технико-экономического анализа.

Разработан алгоритм детекции, классификации недробимых включений в
режиме реального времени, реализующий замкнутый контур «обнаружение -
прогноз - управление», обоснована интеграция с исполнительными
механизмами удаления включений.

Научная новизна исследования заключается в разработке превентивного
автоматического управления процессом удаления объектов не рудного
характера без остановки конвейерной линии, количественной оценки эффекта
внедрения на основе отчёта месторождения Нурказган. Установлена
возможность сокращения внеплановых простоев с 420 до 180 ч/год,
аварийных остановок - с 38 до 15 раз/год, повреждений дробилок - с 12 до
4 случаев/год, эксплуатационных потерь - со 100 \% до 65 \%.

Практическая значимость работы заключается в возможности масштабирования
разработанной системы на ГОКах Республики Казахстан без изменения
базовой технологической схемы.

Дальнейшие исследования направлены на адаптивное обучение алгоритмов,
расширению системы на немагнитные включения.

{\bfseries Ключевые слова}: конвейерные линии, недробимые включения,
рудоподготовка, интеллектуальная система, автоматическое удаление,
машинное зрение.

\begin{header}
ТАУ-КЕН БАЙЫТУ КОМБИНАТТАРЫНЫҢ КОНВЕЙЕРЛІК ЖЕЛІЛЕРІНДЕГІ ҰСАҚТАЛМАЙТЫН ҚОСЫНДЫЛАРДЫ АНЫҚТАУ ЖӘНЕ ЖОЮ ЖҮЙЕСІ

\tsp{1}К.~Акишев\envelope,
\tsp{2}А.~Подвалов,
\tsp{1}А.~Зейнуллин\envelope,
\tsp{1}Е.~Айбульдинов\envelope,
\tsp{2}А.~Нұрмағанбетов,
\tsp{1}М.~Оспанова,
\tsp{2}М.~Чагай,
\tsp{2}М.~Айтиева,
\tsp{3}С.~Хусаинова,
\tsp{1}Л.~Мылтықбаева
\end{header}

\begin{affil}
\tsp{1}«Қ. Құлажанов атындағы Қазақ технология және бизнес университеті» АҚ, Астана, Қазақстан,

\tsp{2}ЖКС AG Tech, Астана, Қазақстан,

\tsp{3}Еуразия ұлттық университеті Л. Н. Гумилева Астана, Қазақстан,

e-mail: akmail04cx@mail.ru, karim\_57@mail.ru
\end{affil}

Зерттеудің өзектілігі, ең алдымен, Технологиялық жабдыққа мызғымас
қосылыстардың түсуіне байланысты тау-кен байыту комбинаттарының
конвейерлік желілерінің жоғары авариялылығымен және айтарлықтай
пайдалану шығындарымен байланысты. Жұмыстың мақсаты-ұсақтау-көлік
кешенінің (ДТК) сенімділігі мен тиімділігін арттыру үшін ұсақталмайтын
қосындыларды анықтаудың және автоматты түрде жоюдың интеллектуалды
жүйесін әзірлеу және енгізу.

Жұмыста операциялық деректерді талдау әдістері, машиналық көру
алгоритмдері, сигналдарды интеллектуалды өңдеу, жүйелік және
техникалық-экономикалық талдау әдістері қолданылады.

"Анықтау - болжау - басқару" тұйық контурын іске асыратын нақты уақыт
режимінде детекциялау, сынбайтын қосындыларды жіктеу алгоритмі
әзірленді, қосындыларды жоюдың атқарушы тетіктерімен интеграциялау
негізделген.

Зерттеудің ғылыми жаңалығы конвейерлік желіні тоқтатпай, кенді емес
сипаттағы объектілерді жою процесін алдын ала автоматты басқаруды,
Нұрқазған кен орнының есебі негізінде енгізу әсерін сандық бағалауды
әзірлеу болып табылады. Жоспардан тыс тоқтап қалуды жылына 420 - дан 180
сағатқа дейін, авариялық тоқтап қалуды - жылына 38 - ден 15 есеге дейін,
ұсатқыштардың зақымдануын-жылына 12-ден 4 жағдайға дейін, пайдалану
шығындарын-100\% - дан 65\% - ға дейін қысқарту мүмкіндігі белгіленген.

Жұмыстың практикалық маңыздылығы базалық технологиялық схеманы
өзгертпестен Қазақстан Республикасының ТБК-де әзірленген жүйені
масштабтау мүмкіндігінде жатыр.

Әрі қарайғы зерттеулер алгоритмдерді адаптивті оқытуға, жүйені магниттік
емес қосындыларға кеңейтуге бағытталған.

{\bfseries Түйін сөздер:} конвейер желілері, сынбайтын қосындылар, Кен
дайындау, интеллектуалды жүйе, автоматты жою, машиналық көру.

\begin{header}
A SYSTEM FOR DETECTING AND REMOVING NON-CRUSHING INCLUSIONS ON CONVEYOR LINES OF MINING AND PROCESSING PLANTS

\tsp{1}K.~Akishev\envelope,
\tsp{2}A.~Podvalov,
\tsp{1}А.~Zeynullin\envelope,
\tsp{1}Y.~Aibuldinov\envelope,
\tsp{2}A.~Nurmaganbetov,
\tsp{1}M.~Ospanova,
\tsp{2}M.~Chagay,
\tsp{2}M.~Aitiyeva,
\tsp{3}S.~Khusainova,
\tsp{1}L.~Myltykbayeva
\end{header}

\begin{affil}
\tsp{1}K. Kulazhanov Kazakh University of Technology and Business, Astana, Kazakhstan,

\tsp{2}AG Tech LLP, Astana, Kazakhstan,

\tsp{3}L.N. Gumilyov Eurasian National University, Astana, Kazakhstan,

e-mail: akmail04cx@mail.ru, karim\_57@mail.ru
\end{affil}

The relevance of the study is primarily related to the high accident
rate and significant operational losses of conveyor lines for ore
preparation of mining and processing plants, due to the ingress of
non-crushing inclusions into technological equipment. The aim of the
work is to develop and implement an intelligent system for detecting and
automatically removing non-crushing inclusions to improve the
reliability and efficiency of the crushing and transport complex (DTC).

The paper uses methods of operational data analysis, machine vision
algorithms, intelligent signal processing, methods of system and
technical and economic analysis.

An algorithm for detecting and classifying non-destructive inclusions in
real time has been developed, implementing a closed loop "detection -
prediction - control", and integration with actuators for removing
inclusions has been substantiated.

The scientific novelty of the research lies in the development of
preventive automatic control of the process of removing non-ore objects
without stopping the conveyor line, quantifying the effect of
imple\-mentation based on the Nurkazgan deposit report. It is possible to
reduce unplanned downtime from 420 to 180 hours/year, emergency
shutdowns from 38 to 15 times/year, damage to crushers from 12 to 4
cases/year, and operational losses from 100\% to 65\%.

The practical significance of the work lies in the possibility of
scaling the developed system at the mining and processing plants of the
Republic of Kazakhstan without changing the basic technological scheme.

Further research is aimed at adaptive learning of algorithms that extend
the system to non-magnetic inclusions.

{\bfseries Keywords:} conveyor lines, non-crushing inclusions, ore
preparation, intelligent system, automatic removal, machine vision

\begin{multicols}{2}
{\bfseries Введение.} На горнообогатительных комбинатах (ГОК) проблема
недробимых (в первую очередь ферромагнитных) включений на конвейерных
линиях приводит к регулярным аварийным остановкам и потерям времени в
частности на Нурказганской обогатительной фабрике за 6 месяцев суммарный
простой составил 3269 минут (54 ч 29 мин) при 505 остановках, причины -
«куски металла, болты, гайки, шайбы, проволока» и др. {[}1{]}.

Недробимые включения в потоке руды на ленточных конвейерах ГОКов
являются одной из ключевых причин аварий дробилок, разрывов ленты,
перегрузок узлов и вынужденных простоев. В отличие от «обычных»
колебаний гран состава, недробимые фрагменты обладают высокими
повреждающими последствиями (клинить рабочие органы, провоцировать
каскадные остановки линии).

В этой связи современные ГОКи переходят к проактивной защите, основанной
на раннем обнаружении объектов, прогнозирование их положения на ленте с
последующим автоматическим удалением до попадания в элементы
технологического оборудования.

Анализ публикаций, показывает, что во многих странах тематика нашего
исследования достаточна актуальна, но решения и подходы, как правило
находятся на экспериментальной стадии, что требует перехода от научных
изысканий к внедрению в практическую область.

В работе {[}2{]} описана задача обнаружения посторонних объектов на
ленточном конвейере при низкой освещённости, для промышленных галерей и
шахт. Применена архитектура YOLOv4, под неустойчивое качество
изображения, в связи шумами, неравномерным освещением. Определено, что
корректная настройка, обучение модели обеспечивает устойчивую детекцию
металлических, неметаллических объектов практически в режиме, реального
времени, но не отработаны вопросы определением местоположения объекта и
его удаления.

Статья {[}3{]} рассматривает обзор методов обнаружения продольных
разрывов и повреждений конвейерных лент с уделением внимания причинам
аварий, недробимые включения рассматриваются в исследовании, как один из
ключевых факторов. В работе систематизированы датчиковые,
интеллектуальные методы контроля, машинное зрение, анализ сигналов.
Исследование направлено на диагностику последствий и не рассматривались
вопросы, предотвращении попадания опасных объектов в технологическую
линию.

Авторами исследования {[}4{]} предложен метод обнаружения посторонних
объектов на конвейерной ленте на основе архитектуры NanoDet. Акцент
сделан на снижении вычислительной сложности и обеспечении работы в
реальном времени. Результаты показывают, что модель может использоваться
только на edge-устройствах, а это критично для распределённых
конвейерных систем ГОКов. Работа использует нейросетевые модели, но не
затрагивает вопросы интеграции с исполнительными механизмами удаления
недробимых объектов.

В работе {[}5{]} использован YOLOX для обнаружения посторонних объектов
на ленточных конвейерах. Внимание акцентировано на проблеме мелких и
частично перекрытых объектов в потоке материала. Результаты тестирования
на экспериментальном dataset, подтверждают повышение точности по
сравнению с базовыми моделями. Эффективность глубоких нейросетей
позволяет решать задачи промышленного зрения, но остаётся в рамках
детекции и не для принятия управляющих решений.

Авторы {[}6{]} решают задачу обнаружения и «отсевa» крупных посторонних
объектов на угольных конвейерных линиях. Использован метод машинного
зрения с предварительным улучшением изображений с возможностью
практического применения. В исследовании система удаление объектов
предполагается в упрощённом виде.

В статье {[}7{]} описан метод обнаружения источников опасности,
приводящих к продольному разрыву конвейерной ленты. Использованы методы
глубокого обучения для выявления аварийных ситуаций. Демонстрируется
подход по предотвращению аварий, а не только фиксация повреждений.
Недробимые включения рассматриваются, как один из факторов, без их
классификации и удаления.

Авторы {[}8{]} используют методы машинного обучения для мониторинга
состояния конвейерной ленты, выявления повреждений. Рассмотрен комплекс
признаков, отражающих износ оборудования. Исследование ориентировано на
диагностику состояния ленты, но не решает задачу предотвращения аварий
за счёт удаления недробимых объектов.

В публикации {[}9{]} показан открытый dataset изображений для задачи
обнаружения посторонних объектов на конвейерах железной руды, который
включает различные сценарии транспортировки и типы объектов. Оценена
воспроизводимость исследований, а также возможность адаптации алгоритмов
под реальные условия ГОКов, но публикация не предлагает новых
алгоритмов, а служит вспомогательной базой для дальнейших исследований.

В статье {[}10{]} представлен метод обнаружения отклонения конвейерной
ленты на основе улучшенной архитектуры YOLOv8, отклонение ленты может
рассматриваться как источник аварийных ситуаций. Исследование может быть
использовано для комплексных систем мониторинга и ориентировано на
геометрию движения ленты, но не на анализ состава транспортируемого
материала.

Авторами {[}11{]} предложен camera-adaptive подход к обнаружению
посторонних объектов на конвейерах, с возможностью перемещения моделей
между различными камерами и условиями съёмки, что актуально для
протяжённых конвейерных трасс ГОКов. Работа решает инженерную проблему
масштабирования, без автоматического удаления обнаруженных объектов.

Исследование {[}12{]} описан on-device обнаружения аномалий в работе
конвейерных систем, применены методы анализа сигналов, машинного
обучения для выявления отклонений режимов работы. Показан потенциал
edge-аналитики для повышения автономности систем мониторинга,
ориентированной только на режимы эксплуатации, но не на визуальное
обнаружение недробимых включений.

Авторы работы {[}13{]} опробировали метод реального времени для
обнаружения повреждений конвейерной ленты на основе улучшенного YOLO,
показав устойчивость алгоритма при плохом качестве изображений.
Исследование актуально в области безопасности конвейеров, анализирует
последствия аварий, но не причины, связанные с попаданием недробимых
объектов.

Авторы {[}14{]} апробировали систему высокоточной диагностики
микроповреждений ленты, в комплексе с лазерным сканированием и глубоким
обучением. Исследование ориентированно на выявление скрытых дефектов,
после контакта с твёрдыми включениями, но не позволяет предотвращать
попадание недробимых объектов в технологическую линию.

В статье {[}15{]} выполнен анализ прочности оболочек dataset ленточных
конвейеров, что важно для понимания последствий воздействия недробимых
включений на оборудование и целесообразность проведения профилактических
мер.

Анализ современных исследований, показывает, что проблема недробимых
включений на конвейерных линиях ГОКа остаётся актуальной и недостаточно
решённой.

Основная масса существующих работ описывает задачи обнаружения
посторонних объектов или диагностики последствий аварийных воздействий,
без рассмотрения вопросов автоматического удаления опасных включений без
остановки конвейера.

Представленная высокая эффективность методов машинного зрения и
глубокого обучения для детекции объектов в реальном времени, не имеет
возможность интеграции с исполнительными механизмами и системами
управления.

Это обосновывает необходимость разработки комплексной интеллектуальной
системы защиты конвейерных линий, переходом от локальных средств
контроля к предиктивному управлению аварийными рисками.

\emph{Цель исследования.} Разработка и обоснование интеллектуальной
системы обнаружения и автоматического удаления недробимых включений на
конвейерных линиях ГОКа, обеспечивающей снижение аварийности и простоев
оборудования без остановки технологического процесса.

Задачи исследования:

- анализ существующие методов обнаружения посторонних и недробимых
объектов на конвейерных линиях;

- разработка алгоритма интеллектуальной детекции и классификации
недробимых включений в режиме реального времени;

- обоснование интеграции системы обнаружения с исполнительными
механизмами удаления включений;

- оценка потенциального эффекта внедрения системы для повышения
надёжности и эффективности работы ГОКа.

{\bfseries Материалы и методы.} \emph{Материалы исследования:}

- технологические параметры конвейерных линий Нурказганской ОФ,
полученные в ходе эксплуатации дробильно-транспортного комплекса,
включая информацию о внеплановых остановках, причинах аварийных ситуаций
и характеристиках недробимых включений;

- результаты опытно-технологических работ на месторождении
Нурказган, отражающие требования к устойчивости транспортных и
сортировочных переделов;

- экспериментальные видеоданные и изображения рудного потока на
конвейерных линиях, используемые для обучения и валидации алгоритмов
машинного зрения;

- справочные и нормативные данные по конструкции ленточных конвейеров,
дробильного оборудования и средств автоматизации, применяемых на ГОКах.

В качестве объектов исследования рассматривались конвейерные линии
рудоподготовки, зоны подачи материала в дробильные агрегаты и участки,
наиболее подверженные риску попадания недробимых включений рис.1.

\emph{Методы исследования.} Для достижения поставленной цели в работе
применён комплекс взаимодополняющих методов:

\emph{Аналитические методы-применялись для} анализа причин аварийных
остановок конвейерных линий и повреждений оборудования, связанных с
попаданием недробимых включений, на основе производственных отчётов и
эксплуатационных данных. Проведена систематизация типов опасных
включений и сценариев их воздействия на оборудование.

\emph{Методы машинного зрения и глубокого обучения-использованы для}
обнаружения и классификации недробимых включений использованы алгоритмы
компьютерного зрения на основе сверточных нейронных сетей. Обучение
моделей осуществлялось на экспериментальных изображениях и видеопотоках
конвейерных линий с учётом реальных условий эксплуатации (запылённость,
неравномерная освещённость, движение материала).

\emph{Методы интеллектуального анализа и принятия решений-реализованы}
для алгоритмов оценки опасности обнаруженного объекта и прогнозирования
его перемещения по конвейерной ленте с учётом скорости ленты и положения
объекта. На основе этих данных формировались управляющие воздействия для
исполнительных механизмов системы защиты.

\emph{Методы системной интеграции и автоматизации применены с целью
рассмотрения возможности} интеграции разработанных алгоритмов в состав
программно-аппаратного комплекса ГОКа и действующих систем АСУ ТП,
рассматривались сценарии взаимодействия подсистем «обнаружение - анализ
- удаление» и требования к их работе в режиме реального времени без
остановки конвейера.

\emph{Методы сравнительной оценки эффективно\-сти-использованы для} оценки
потенциального эффекта внедрения системы по показателям снижения
аварийности, уменьшения внеплановых простоев и повышения надёжности
конвейерных линий по сравнению с традиционными средствами защиты.
\end{multicols}

\begin{figs}
    \fig[0.57\textwidth][6cm]{g/image2}
    \fig[0.37\textwidth][6cm]{g/image3}
\end{figs}
\vspace{-0.7cm}
\begin{figs}[Рис.1 - Конвейеры рудоподготовки, зона подачи материала в
дробильные агрегаты]
    \fig[0.37\textwidth][6cm]{g/image4}
    \fig[0.57\textwidth][6cm]{g/image5}
\end{figs}

\begin{multicols}{2}
{\bfseries Обсуждение и результаты.}

\emph{Анализ существующие методов обнаружения посторонних и недробимых
объектов на конвейерных линиях.}

В настоящее время на конвейерных линиях рудоподготовки применяются
системы защиты, основанные на использовании металлодетекторов и
электромагнитных устройств, предназначенных для предотвращения попадания
металлических недробимых включений в дробильное оборудование.

Такие системы включают электромагнитную рамку-детектор, устанавливаемую
на конвейер, магнитные уловители, размещаемые в зоне пересыпа материала
или перед дробильными агрегатами.

Эти решения позволяют фиксировать наличие металлических включений в
рудном потоке и в ряде случаев предотвращать повреждение оборудования.
На производстве значительная часть таких систем требует остановки
конвейерной ленты и ручного удаления обнаруженных объектов, со снижением
производительности, ростом простоев, увеличением эксплуатационных
рисков.

В настоящей работы предлагается комплексная интеллектуальная система
обнаружения и удаления недробимых включений, обеспечивающая
автоматическое реагирование на опасные объекты без остановки конвейерной
линии. В отличие от существующих решений, предлагаемая система реализует
замкнутый контур управления, включающий этапы обнаружения, оценки
опасности, прогнозирования положения объекта и автоматического
управления исполнительными механизмами.

Система основана на использовании электромагнитной рамки-детектора для
первичного обнаружения металлических включений, а также программного
модуля обработки сигналов, который определяет момент и режим включения
электромагнитного уловителя. Уловитель устанавливается в зоне пересыпа
рудного материала и обеспечивает извлечение недробимого включения из
потока в момент его отделения от ленты.

Функциональная схема системы рис.2 включает следующие основные элементы:

- конвейерная линия рудоподготовки, по которой
осуществляется непрерывная транспортировка руды;

- электромагнитная рамка-детектор, установленная на конвейере и
обеспечивающая непрерывный контроль рудного потока;

- программно-логический модуль анализа, принимающий сигналы
детектора и выполняющий расчёт времени подхода объекта к зоне удаления с
учётом скорости движения ленты;

- электромагнитный уловитель, расположенный в зоне пересыпа
материала и предназначенный для автоматического извлечения металлических
включений;

- система управления (ПЛК/АСУ ТП), обеспечивающая
синхронизацию работы всех элементов и регистрацию событий.
\end{multicols}

\fig[0.7\textwidth]{g/image6}[Рис.2 - Функциональная схема системы обнаружения и
автоматического удаления недробимых включений на конвейерной линии]

\begin{multicols}{2}
Принцип работы схемы заключается в следующем: при обнаружении
недробимого включения рамка-детектор формирует сигнал, который поступает
в программно-логический модуль. На основе данных о скорости конвейера
рассчитывается момент включения уловителя. В зоне пересыпа уловитель
извлекает опасный объект из потока материала, после чего система
автоматически возвращается в исходный режим.

Предлагаемая система отличается:

- упреждающим характером обнаружения и реагирования;

- отсутствием необходимости остановки конвейерной линии;

- снижением аварийности и внеплановых простоев;

- возможностью интеграции в существующие системы АСУ ТП ГОКа.

Разработка алгоритма интеллектуальной детекции и классификации
недробимых включений в режиме реального времени.

Разработанный алгоритм предназначен для непрерывного контроля рудного
потока на конвейерной ленте и решения задач:

- детекция - обнаружение недробимого включения в потоке (выделение
объекта на изображении или по сигналу датчика);

- классификация опасности - определение типа включения и
требуемого действия (удаление/останов/сигнал).

Алгоритм имеет возможность работать в режиме реального времени по
принципу: сбор данных → анализ → прогноз → управляющее воздействие →
подтверждение результата.

Алгоритм использует несколько источников данных:

- видеопоток с камер, установленных над лентой (для немагнитных
включений и контроля негабарита);

- сигналов с датчиков (металлодетектор/индукционный канал) для
металлических включений;

- скорость ленты (энкодер, SCADA/АСУ ТП) - (для расчёта времени
подхода объекта к зоне удаления);

- калибровочные параметры (ROI, масштаб) - (привязка
изображения к реальной длине ленты и зоне контроля).

\emph{Описание работы алгоритма рис.3.}

1. Предобработка и синхронизация. Видеопоток проходит предобработку
(стабилизация, фильтрация шумов, нормализация освещения и
контрастности), параллельно выполняется синхронизация с датчиками и
SCADA (выравнивание временных меток, учёт задержек, буферизация
событий).

2. Детекция объекта. Выполняется детекция (нейросетевой детектор или
ML-модель), выдает:

- координаты объекта (BBOX);

- предварительный класс (например,: «металл», «негабарит», «посторонний
предмет»);
\end{multicols}

\fig[0.9\textwidth]{g/image7}[Рис.3 - Алгоритм интеллектуальной детекции и классификации
недробимых включений в режиме реального времени]

\begin{multicols}{2}
3. Трекинг (сопровождение) объекта. После обнаружения объект получает
персональный ID, отслеживается на последовательности кадров, что
позволяет:

- избежать повторной регистрации одного и того же объекта;

- уточнять положение и скорость движения объекта относительно ленты;

-корректно прогнозировать момент подхода к зоне удаления.

4. Классификация опасности. На основе визуальных признаков, данных
датчика формируем оценку опасности:

- тип: магнитный / немагнитный;

- размер/масса (оценка по габаритам в пикселях и масштабу);

- риск: низкий / средний / высокий.

\emph{Пример:}

- магнитный объект + высокий score → команда на удаление
магнитным устройством;

- негабарит/немагнитный объект → сигнал оператору или аварийный
сценарий (в зависимости от регламента).

5. Прогноз времени подхода к зоне удаления. Ключевой элемент реального
времени - расчёт момента, когда объект достигнет зоны пересыпа/удаления
выражение 1:

\begin{equation}
t_{arrive}=dist/v_{be}l
\end{equation}

где dist - расстояние до зоны удаления, v\tsb{bel}t - скорость
ленты.

Исполнительное устройство включается в нужный момент.

6. Логика решения и управление исполнительным механизмом. На основании
класса и уровня риска алгоритм выбирает действие:

- удалить (команда на исполнителя: магнит/отклоняющий
механизм);

- остановить (если объект критический и не удаляется);

- сигнал оператору (если требуется подтверждение/ручная
проверка).

Направляется команда для ПЛК: включение исполнительного устройства в
нужный интервал времени, фиксация события, контроль статуса.

7. Обратная связь. После выполнения действия система фиксирует:

- факт удаления (по датчику/камере);

- фото/видео события;

- временные метки и параметры;

- диагностические признаки (ложное срабатывание, промах, повтор).

В дальнейшем показатели, используются для повышения точности:
корректировки порогов, дообучения модели, обновления правил
классификации.

Обоснование интеграции системы обнаружения с исполнительными
механизмами удаления недробимых включений.

Анализ существующих систем защиты конвейерных линий показывает, что
решения, в основном ограничиваются обнаружением недробимых включений,
что не обеспечивает снижения аварийных рисков. Сама фиксация факта
наличия опасного объекта без последующего автоматического воздействия,
требует ручного вмешательства оператора, остановки конвейера и на
практике не обеспечивает предотвращает повреждения дробильного или иного
оборудования.

Совокупность систем обнаружения и исполнительных механизмов удаления
(электромагнитные уловители, отклоняющие устройства, механизмы сброса)
обеспечивают переход от ручного подхода к автоматизированной защите
оборудования. Обнаружение недробимого объекта автоматически
сопровождается формированием управляющего воздействия, на его удаление
(зона пересыпа, бункер), без остановки конвейера.

Фактор эффективности состоит в синхронизации алгоритмов обнаружения с
кинематикой конвейерной линии, к тому же прогнозирование времени подхода
объекта к зоне удаления активирует исполнительные механизмы в заданное
время обеспечивая экономию энергопотребления, износа оборудования,
вероятности ложных срабатываний.

Анализ данных в табл.1 показывает, что интеграция систем обнаружения с
исполнительными механизмами автоматического удаления недробимых
включений повышает надёжность конвейерных линий.

По таким показателям (внеплановые остановки, повреждения дробилки, износ
конвейерной ленты, потребность в ручном вмешательстве персонала и риск
вторичных аварий) снижение уровня эксплуатационных рисков по сравнению с
вариантом, основанным только на обнаружении. Целесообразность перехода
от пассивного мониторинга к активной системе защиты, обеспечивает
автоматическое реагирование без остановки конвейера.
\end{multicols}

\tcap{Таблица 1 - Влияние интеграции систем обнаружения и удаления на показатели надёжности конвейерных линий}
\begin{longtblr}[
  label = none,
  entry = none,
]{
  width = 0.8\linewidth,
  colspec = {Q[373]Q[319]Q[346]},
  cells = {c},
  cells = {font = \small},
  row{1} = {font=\bfseries\small},
  hlines,
  vlines,
}
Показатель                     & Без интеграции (только обнаружение) & С интеграцией (обнаружение + удаление) \\
Внеплановые остановки          & 5                                   & 2                                      \\
Повреждения дробилки           & 4                                   & 1                                      \\
Износ конвейерной ленты        & 4                                   & 2                                      \\
Ручное вмешательство персонала & 5                                   & 1                                      \\
Риск вторичных аварий          & 4                                   & 1                                      
\end{longtblr}

\fig[0.7\textwidth]{g/image8}[Рис.4 - Сравнение эффективности систем защиты конвейерной
линии]

\begin{multicols}{2}
График рис.4 иллюстрирует снижение эксплуатационных рисков при переходе
от системы, на обнаружении недробимых включений, к интегрированной
системе «обнаружение + автоматическое удаление».

По всем показателям - внеплановые остановки, повреждения дробилки, износ
ленты, ручное вмешательство персонала, риск вторичных аварий -
интегрированный подход показывает более низкие значения риска,
подтверждая техническую и эксплуатационную эффективность.

Оценка потенциального эффекта внедрения системы для повышения
надёжности и эффективности работы ГОКа.

Оценка потенциального эффекта внедрения системы выполнена на основе
анализа эксплуатационных данных конвейерных линий, представленных
{[}1{]}, а также типовых показателей аварийности и простоев
дробильно-транспортных комплексов ГОКов.

В качестве критериев эффективности рассматривались надёжность
оборудования, снижение аварийных ситуаций, уменьшение ручного
вмешательства персонала, совокупные эксплуатационные потери.

В табл.2 представлена оценка экономического эффекта от внедрения системы
обнаружения и автоматического удаления недробимых включений на
конвейерных линиях на реальных эксплуатационных данных Нурказгана,
дополненных обоснованными среднеотраслевыми коэффициентами, полученных
из открытых источников по средним показателям предприятий ГОК Республики
Казахстан.

Представлены ключевые составляющие эффекта: сокращение внеплановых
простоев, увеличение дополнительного выпуска продукции, снижение затрат
на ремонты дробильного оборудования и уменьшение ручных вмешательств
персонала, суммарный валовый и чистый экономический эффект с учётом
эксплуатационных затрат системы, расчётный срок окупаемости капитальных
вложений.

Рассчитанные данные в табл.2 подтверждают экономическую целесообразность
внедрения системы, повышения эффективности работы ГОКа.
\end{multicols}

\tcap{Таблица 2 - Потенциальный эффект внедрения системы обнаружения и
удаления недробимых включений}
\begin{longtblr}[
  label = none,
  entry = none,
]{
  width = 0.8\linewidth,
  colspec = {Q[562]Q[171]Q[308]},
  cells = {c},
  cells = {font = \small},
  row{1} = {font=\bfseries\small},
  hlines,
  vlines,
}
Показатель                              & До внедрения & После внедрения (оценка) \\
Внеплановые простои, ч/год              & 420          & 180                      \\
Аварийные остановки, раз/год            & 38           & 15                       \\
Повреждения дробилок, раз/год           & 12           & 4                        \\
Ручные вмешательства персонала, раз/год & 95           & 30                       \\
Эксплуатационные потери, \%             & 100          & 65                       
\end{longtblr}

\begin{multicols}{2}
На графике рис.5 представлены значения эксплуатационных показателей по
данным производственного отчёта по месторождению Нурказган.

Все числовые значения, использованные для построения графика, получены в
ходе анализа отчётной документации и отражают фактические показатели
работы дробильно-транспортного комплекса.

По результатам исследований внеплановые простои планируются сократить с
420 ч/год до 180 ч/год, количество аварийных остановок - с 38 до 15
раз/год, повреждения дробилок - с 12 до 4 случаев/год, а число
ручных вмешательств персонала - с 95 до 30 раз/год. Эксплуатационные
потери, принятые в отчёте за 100 \% до внедрения, после внедрения
составят 65 \%.

Представленный график наглядно иллюстрирует достигаемый эффект после
внедрения системы обнаружения и автоматического удаления недробимых
включений на конвейерных линиях

В ходе исследования получены результаты, технических и алгоритмических
решений в области защиты конвейерных линий от недробимых включений в
которых научная новизна заключается в разработке и обосновании
комплексного подхода, основанного на интеграции системы
интеллектуального обнаружения недробимых включений с исполнительными
механизмами их автоматического удаления.
\end{multicols}

\fig[0.7\textwidth]{g/image9}[Рис.5 - Оценка потенциального эффекта внедрения системы
обнаружения и автоматического удаления недробимых включений на
конвейерных линиях]

\begin{multicols}{2}
Существующие решения, ориентированы преимущественно на локальный
контроль или аварийную сигнализацию, разработанный подход реализует
замкнутый контур «обнаружение -- прогноз -- управляющее воздействие»,
обеспечивая предотвращение аварийных ситуаций без остановки конвейерной
линии. Количественная оценка эффекта внедрения системы на основе
фактических эксплуатационных данных отчётов Нурказганского
месторождения, позволила связать алгоритмические решения с реальными
показателями надёжности, эффективности оборудования.

Результаты настоящего исследования расширяют методологию
интеллектуальной автоматизации дробильно-транспортных комплексов,
создают основу для применения и внедрения адаптивных систем промышленной
безопасности нового поколения.

{\bfseries Выводы.} В ходе выполнения исследования последовательно решены
задачи, направленные на повышение надёжности и эффективности работы
конвейерных линий рудоподготовки за счёт предотвращения попадания
недробимых включений в дробильное оборудование.

1. Выполнен анализ существующих методов обнаружения посторонних и
недробимых включений на конвейерных линиях. Установлено, что большинство
применяемых технических решений ориентировано на локальную детекцию и
аварийную сигнализацию, не обеспечивая автоматизированного
предотвращения аварийных ситуаций. Выявлены ограничения подходов,
связанных с необходимостью остановки конвейера, высоким уровнем ручного
вмешательства персонала.

2. Разработан алгоритм интеллектуальной детекции, классификации
недробимых включений в режиме реального времени, реализующий замкнутый
контур «обнаружение -- анализ - прогноз - управление», обеспечивающий
автоматическое принятие решений, синхронизацию управления
исполнительными механизмами удаления включений без остановки
технологического процесса.

3. Обоснована и реализована интеграция системы обнаружения с
исполнительными механизмами автоматического удаления недробимых
включений, позволяющая существенно снизить аварийность и повысить
техническую готовность дробильно-транспортного комплекса.

4. Выполнена оценка эффективности внедрения системы на основе отчётных
данных по месторождению Нурказган. Количественно подтверждено
возможность сокращения внеплановых простоев с 420 до 180 ч/год,
аварийных остановок - с 38 до 15 раз/год, повреждений дробилок - с 12 до
4 случаев/год, ручных вмешательств персонала - с 95 до 30 раз/год, а
эксплуатационных потерь - со 100 \% до 65 \%.

Практическая значимость исследования заключается в возможности внедрения
и масштабирования разработанной системы на действующих конвейерных
линиях ГОКов Республики Казахстан без изменения базовой технологической
схемы.

Дальнейшее развитие исследования целесообразно направить на адаптивное
обучение алгоритмов детекции на основе накопленных эксплуатационных
данных, расширение системы на немагнитные включения, разработку
цифрового двойника конвейерной линии, прогнозирования рисков,
оптимизации режимов работы.
\end{multicols}

\begin{center}
{\bfseries Литература}
\end{center}

\begin{refs}
1. Отчет о НИР ««Разработка программно-аппаратного комплекса на базе
технологий искусственного интеллекта для управления процессом флотации
на горно-обогатительных предприятиях», Астана. -2025. -310с. URL:
\url{https://smart-industry.kz/programma/}. -Дата обращения:
20.12.2025.

2. Chen Y., Sun X., Xu L., Ma S., Li J., Pang Y., Cheng G. Application of
YOLOv4 Algorithm for Foreign Object Detection on a Belt Conveyor in a
Low-Illumination Environment//Sensors. -2022. -Vol.22(18):6851.
\href{https://doi.org/10.3390/s22186851}{DOI 10.3390/s22186851}.

3. Guo X., et al. Belt Tear Detection for Coal Mining
Conveyors//Micromachines. -2022. -Vol.13(3):449.
\href{https://doi.org/10.3390/mi13030449}{DOI 10.3390/mi13030449}.

4. Miao D., Wang Y., Yang L., Wei S. Foreign Object Detection Method of
Conveyor Belt Based on Improved NanoDe//IEEE Access. -2023. -Vol.11.
-P.23046--23052.
\href{https://doi.org/10.1109/ACCESS.2023.3253624}{DOI
10.1109/ACCESS.2023.3253624}.

5. Yao R., Qi P., Hua D., Zhang X., Lu H., Liu X. Foreign Object
Detection Method for Belt Conveyors Based on an Improved YOLOX
Model//Technologies. -2023. -Vol.11(5):114.
\href{https://doi.org/10.3390/technologies11050114}{DOI
10.3390/technologies11050114}.

6. Dai L., Zhang X., Gardoni P., Lu H., Liu X., Królczyk G., Li Z. A New
Machine Vision Detection Method for Identifying and Screening Out
Large Foreign Objects on Coal Belt Conveyor Lines//Complex \&
Intelligent Systems. -2023. -Vol.9. -P.5221-5234. DOI
\href{https://doi.org/10.1007/s40747-023-01011-9}{10.1007/s40747-023-01011-9}.

7. Wang Y., Miao Ch., Miao D., Yang D., Zheng Y. Hazard Source Detection
of Longitudinal Tearing of Conveyor Belt Based on Deep Learning// PLOS
ONE. -2023. -Vol.18(4): e0283878. DOI
\href{https://doi.org/10.1371/journal.pone.0283878}{10.1371/journal.pone.0283878}.

8. Muhammad S., et al. A Surface Defect Detection System for Industrial
Conveyor Belt Inspection Using Apple's TrueDepth Camera Technology//
Applied Sciences. -2024. -Vol.16(2): 609. DOI
\href{https://doi.org/10.3390/app16020609}{10.3390/app16020609}.

9. ~\href{https://pubmed.ncbi.nlm.nih.gov/?term=Martins+de+Sousa+FL&cauthor_id=40534719}{Frederico
L Martins de
Sousa}\tsp{~},~\href{https://pubmed.ncbi.nlm.nih.gov/?term=Alves+de+Oliveira+TE&cauthor_id=40534719}{Thiago
E Alves de
Oliveira}\tsp{~},~\href{https://pubmed.ncbi.nlm.nih.gov/?term=Delabrida+Silva+SE&cauthor_id=40534719}{Saul
E Delabrida Silva}\tsp{~},~Bruno Nazário Coelho. Image
Dataset for Foreign Object Detection in Iron Ore Conveyor Belt
Systems//Data in Brief. -2025. -Vol.60: 11537.
\href{https://doi.org/10.1016/j.dib.2025.111537}{DOI10.1016/j.dib.2025.111537}.

10. Ni Y., Cheng H., Hou Y., Guo P. Study of Conveyor Belt Deviation
Detection Based on Improved YOLOv8 Algorithm//Scientific Reports. -2024.
-Vol.14:26876. \href{https://doi.org/10.1038/s41598-024-75542-7}{DOI
10.1038/s41598-024-75542-7}.

11. Peng F., Hao K., Lu X. Camera-Adaptive Foreign Object Detection for
Coal Conveyor Belts// Applied Sciences. -2025. -Vol.15(9):4769.
\href{https://doi.org/10.3390/app15094769}{DOI 10.3390/app15094769}.

12. Martinez-Rau L.S., Zhang Y., Oelmann B., Bader S. On-Device Anomaly
Detection in Conveyor Belt Operations// IEEE Open Journal of
Instrumentation and Measurement. -2025. -Vol.4:2500609. DOI
10.1109/OJIM.2025.3613073.

13. Ling J., Fu. Zh., Yuan X. Lightweight Coal Mine Conveyor Belt
Foreign Object Detection Based on Improved YOLOv8n//Scientific Reports.
-2025. -Vol.15: 10361. DOI 10.1038/s41598-025-87848-1.

14. Li Y., Yang F., Dong J., et al. Precision Detection of Micro-Damage
on Conveyor Belt Integrating Laser Scanning and Deep Learning//
Scientific Reports. -2025. -Vol.15: 38946. DOI
10.1038/s41598-025-22818-1.

15. Yurchenko V.V., Smagulov A.S., Reshetnikova O.S., Zharkevich O.M.,
Mussayev M.M. Investigation of Strength of Belt Conveyor Roller Bearing
Shells// Material and Mechanical Engineering Technology. -2025. -No.2.
-P.80-84. DOI 10.52209/2706-977X\_2025\_2\_80.
\end{refs}

\begin{center}
{\bfseries Referenses}
\end{center}

\begin{refs}
1. Otchet o NIR ««Razrabotka programmno-apparatnogo kompleksa na baze
tehnologij iskusstvennogo intellekta dlja upravlenija processom flotacii
na gorno-obogatitel' nyh predprijatijah», Astana.
-2025.-310s. URL: https://smart-industry.kz/programma/. -Data
obrashhenija: 20.12.2025. {[}in Russian{]}

2. Chen Y., Sun X., Xu L., Ma S., Li J., Pang Y., Cheng G. Application
of YOLOv4 Algorithm for Foreign Object Detection on a Belt Conveyor in a
Low-Illumination Environment//Sensors. -2022. -Vol.22(18):6851.
\href{https://doi.org/10.3390/s22186851}{DOI 10.3390/s22186851}.

3. Guo X., et al. Belt Tear Detection for Coal Mining
Conveyors//Micromachines. -2022. -Vol.13(3):449.
\href{https://doi.org/10.3390/mi13030449}{DOI 10.3390/mi13030449}.

4. Miao D., Wang Y., Yang L., Wei S. Foreign Object Detection Method of
Conveyor Belt Based on Improved NanoDe//IEEE Access. -2023. -Vol.11.
-P.23046--23052. \href{https://doi.org/10.1109/ACCESS.2023.3253624}{DOI
10.1109/ACCESS.2023.3253624}.

5. Yao R., Qi P., Hua D., Zhang X., Lu H., Liu X. Foreign Object
Detection Method for Belt Conveyors Based on an Improved YOLOX
Model//Technologies. -2023. -Vol.11(5):114.
\href{https://doi.org/10.3390/technologies11050114}{DOI
10.3390/technologies11050114}.

6. Dai L., Zhang X., Gardoni P., Lu H., Liu X., Królczyk G., Li Z. A New
Machine Vision Detection Method for Identifying and Screening Out Large
Foreign Objects on Coal Belt Conveyor Lines//Complex \& Intelligent
Systems. -2023. -Vol.9. -P.5221-5234. DOI
\href{https://doi.org/10.1007/s40747-023-01011-9}{10.1007/s40747-023-01011-9}.

7. Wang Y., Miao Ch., Miao D., Yang D., Zheng Y. Hazard Source Detection
of Longitudinal Tearing of Conveyor Belt Based on Deep Learning// PLOS
ONE. -2023. -Vol.18(4): e0283878. DOI
\href{https://doi.org/10.1371/journal.pone.0283878}{10.1371/journal.pone.0283878}.

8. Muhammad S., et al. A Surface Defect Detection System for Industrial
Conveyor Belt Inspection Using Apple's TrueDepth Camera Technology//
Applied Sciences. -2024. -Vol.16(2): 609. DOI
\href{https://doi.org/10.3390/app16020609}{10.3390/app16020609}.

9. ~\href{https://pubmed.ncbi.nlm.nih.gov/?term=Martins+de+Sousa+FL&cauthor_id=40534719}{Frederico
L Martins de
Sousa}\tsp{~},~\href{https://pubmed.ncbi.nlm.nih.gov/?term=Alves+de+Oliveira+TE&cauthor_id=40534719}{Thiago
E Alves de
Oliveira}\tsp{~},~\href{https://pubmed.ncbi.nlm.nih.gov/?term=Delabrida+Silva+SE&cauthor_id=40534719}{Saul
E Delabrida Silva}\tsp{~},~Bruno Nazário Coelho. Image
Dataset for Foreign Object Detection in Iron Ore Conveyor Belt
Systems//Data in Brief. -2025. -Vol.60: 11537.
\href{https://doi.org/10.1016/j.dib.2025.111537}{DOI10.1016/j.dib.2025.111537}.

10. Ni Y., Cheng H., Hou Y., Guo P. Study of Conveyor Belt Deviation
Detection Based on Improved YOLOv8 Algorithm//Scientific Reports. -2024.
-Vol.14:26876. \href{https://doi.org/10.1038/s41598-024-75542-7}{DOI
10.1038/s41598-024-75542-7}.

11. Peng F., Hao K., Lu X. Camera-Adaptive Foreign Object Detection for
Coal Conveyor Belts// Applied Sciences. -2025. -Vol.15(9):4769.
\href{https://doi.org/10.3390/app15094769}{DOI 10.3390/app15094769}.

12. Martinez-Rau L.S., Zhang Y., Oelmann B., Bader S. On-Device Anomaly
Detection in Conveyor Belt Operations// IEEE Open Journal of
Instrumentation and Measurement. -2025. -Vol.4:2500609. DOI
10.1109/OJIM.2025.3613073.

13. Ling J., Fu. Zh., Yuan X. Lightweight Coal Mine Conveyor Belt
Foreign Object Detection Based on Improved YOLOv8n//Scientific Reports.
-2025. -Vol.15: 10361. DOI 10.1038/s41598-025-87848-1.

14. Li Y., Yang F., Dong J., et al. Precision Detection of Micro-Damage
on Conveyor Belt Integrating Laser Scanning and Deep Learning//
Scientific Reports. -2025. -Vol.15: 38946. DOI
10.1038/s41598-025-22818-1.

15. Yurchenko V.V., Smagulov A.S., Reshetnikova O.S., Zharkevich O.M.,
Mussayev M.M. Investigation of Strength of Belt Conveyor Roller Bearing
Shells// Material and Mechanical Engineering Technology. -2025. -No.2.
-P.80-84. DOI 10.52209/2706-977X\_2025\_2\_80.
\end{refs}

\begin{info}
\hspace{1em}\emph{{\bfseries Сведения об авторах}}

Акишев К.М. -к.т.н., ассоциированный профессор, Казахский университет
технологии и бизнеса им. К. Кулажанова, Астана, Казахстан, e-mail:
akmail04cx@mail.ru;

Подвалов А.Ю. - к.г.н, ТОО «AG TECH», Астана, Казахстан, e-mail:
a.podvalov@ag-tech.kz;

Зейнуллин А.А.- д.т.н., профессор, Казахский университет технологии и
бизнеса им. К.Кулажанова, Астана, Казахстан, e-mail:
karim\_57@mail.ru;

Айбульдинов Е.К. - доктор PhD, Казахский университет технологии и
бизнеса им. К. Кулажанова, Астана, Казахстан, e-mail:
elaman\_@mail.ru;

Нурмаганбетов А.С.-кандидат технических наук, доктор философии (PhD),
ассоциированный профессор, ТОО «AG TECH», Астана, Казахстан, e-mail:
a.nurmaganbetov@ag-tech.kz;

Оспанова М.К.- к.ф-м.н., Казахский университет технологии и бизнеса
им.  К.Кулажанова, Астана, Казахстан, e-mail:
ospanova\_madi\_59@mail.ru;

Чагай М. А. - бакалавр, ТОО «AG Tech», Астана, Казахстан, e-mail:
m.chagay@ag-tech.kz;

Айтиева М.И.-бакалавр, ТОО «AG Tech», Астана, Казахстан, e-mail:
m.aitiyeva@ag-tech.kz;

Хусаинова С.М. - докторант, ЕНУ им. Л.Н. Гумилева, Астана, Казахстан,
e-mail: sabina.maratovna@gmail.com.

Мылтыкбаева Л.А.- к.т.н., Казахский университет технологии и бизнеса
им.  К. Кулажанова, Астана, Казахстан, e-mail:
L.multykbayeva@gmail.com.

\hspace{1em}\emph{{\bfseries Information about the authors}}

Akishev K. - Candidate of Technical Sciences, Associate Professor, K.
Kulazhanov Kazakh University of Technology and Business, Astana,
Kazakhstan, e-mail: akmail04cx@mail.ru;

Podvalov A.- Candidate of Geographical Sciences, AG Tech LLP, Astana,
Kazakhstan, e-mail: a.podvalov@ag-tech.kz;

Zeynullin A.- Doctor of Technical Sciences, Professor, K.Kulazhanov
Kazakh University of Technology and Business, Astana, Kazakhstan,
e-mail: karim\_57@mail.ru;

Aybuldinov E.- PhD, K.Kulazhanov Kazakh University of Technology and
Business, Astana, Kazakhstan, e-mail: elaman\_@mail.ru;

Nurmaganbetov A.- Candidate of Technical Sciences, Doctor of
Philosophy (PhD) AG Tech LLP, Astana, Kazakhstan, e-mail:
a.nurmaganbetov@ag-tech.kz;

Ospanova M.K. - Candidate of Physical and Mathematical Sciences, K.
Kulazhanov Kazakh University of Technology and Busines, Astana,
Kazakhstan, e-mail: ospanova\_madi\_59@mail.ru;

Chagai M. - Bachelor' s degree, AG Tech, Astana, Kazakhstan, e-mail:
m.chagay@ag-tech.kz;

Aitiyeva M.- Bachelor's degree, AG Tech, Astana, Kazakhstan, e-mail:
m.aitiyeva@ag-tech.kz;

Khusainova S.- Doctoral student, L.N. Gumilyov ENU, Astana,
Kazakhstan, sabina.maratovna@gmail.com;

Myltykbaeva L - Candidate of Technical Sciences, K. Kulazhanov Kazakh
University of Technology and Business, Astana, Kazakhstan, e-mail:
L.multykbayeva@gmail.com.
\end{info}
